% part: intuitionistic-logic
% chapter: semantics
% section: introduction

\documentclass[../../../include/open-logic-section]{subfiles}

\begin{document}

\olfileid{int}{sem}{int}

\olsection{سريزه}

له معنايي پوهنې پرته د هېڅ منطق بشپړ قناعت وړ تشريح نۀ کېږي، او
شهودي منطق هم استثنا نۀ دے۔ د کلاسيک منطق لپاره په !!{valuation}s
ولاړه معنايي پوهنه معياري ده، خو د شهودي منطق لپاره څو بېلابېلې
معنايي پوهنې سيالي کوي۔ هېڅ يوه ئې په دې معنا بشپړه قناعت وړ نۀ ده
چې د نښلوونکو د معناوو د شهودي منلو وړ تشريح ورکړي۔

د !!{relational model}s پر بنسټ معنايي پوهنه، چې د وجهي منطقونو له
معنايي پوهنې سره ورته ده، ښايي تر ټولو ډېره دود وي۔ په دې معنايي
پوهنه کښې !!{propositional variable}s نړۍو ته ټاکل کېږي او دغه
نړۍې د لاسرسي په اړيکه تړلې وي۔ دا اړيکه تل جزئي ترتيب وي، يعنې
انعکاسي، ضد متقارن او متعدي وي۔

په شهودي ډول دا نړۍې د پوهې حالتونه يا «د شواهدو حالتونه» ګڼلے شئ۔
حالت~$w'$ له~$w$ څخه هله او يوازې هله لاسرسي وړ دے چې زموږ د
اوسنۍ پوهې له مخې $w'$ د پوهې يو ممکن راتلونکے حالت وي؛ يعنې له
هغه څۀ سره سمون ولري چې په~$w$ کښې معلوم دي۔ کله چې يوه قضيه
معلومه شي، بيا نامعلومه کېدے نۀ شي: که $!A$ په~$w$ کښې معلوم وي
او~$Rww'$ وي، نو $!A$ په~$w'$ کښې هم معلوم وي۔ نو «پوهه» د
لاسرسي د اړيکې له مخې يورخيزه زياتېدونکې ده۔

که د معرفتي منطق په څېر «$!A$ معلوم دے» داسې تعريف کړو چې «په ټولو
معرفتي بديلو کښې رښتيا دے»، نو $!A \land !B$ په~$w$ کښې هله
معلوم وي چې په ټولو معرفتي بديلو کښې $!A$ او~$!B$ دواړه معلوم وي۔
خو څرنګه چې پوهه يورخيزه زياتېږي او $R$ انعکاسي ده، دا معنا لري
چې $!A \land !B$ په $w$ کښې هله او يوازې هله معلوم وي چې $!A$
او $!B$ په~$w$ کښې دواړه معلوم وي۔ د همدې دليل له مخې
$!A \lor !B$ په $w$ کښې هله او يوازې هله معلوم وي چې لږ تر لږه
يو غړی ئې معلوم وي۔ نو د $\land$ او $\lor$ د رښتياوالي شرطونه
د کلاسيک منطق له شرطونو سره يو شان دي۔

خو د شرطي د رښتياوالي شرطونه له کلاسيک منطق سره توپير لري۔
$!A \lif !B$ په~$w$ کښې هله او يوازې هله معلوم وي چې داسې هېڅ
$w'$، چې $Rww'$ ولري، نۀ وي چې هلته $!A$ معلوم وي خو $!B$
معلوم نۀ وي۔ دا له دې شرط سره يو شان نۀ دے چې $!A$ نامعلوم وي
يا $!B$ په~$w$ کښې معلوم وي۔ ځکه که په~$w$ کښې نۀ $!A$
او نۀ $!B$ معلوم وي، ښايي داسې راتلونکے معرفتي حالت~$w'$ وي
چې $Rww'$ ولري او په $w'$ کښې $!A$ معلوم شي خو~$!B$ معلوم نۀ شي۔

موږ $\lnot !A$ يوازې هغه وخت معلومولے شو چې داسې هېڅ ممکن
راتلونکے معرفتي حالت نۀ وي چې پکې~$!A$ معلوم وي۔ دلته تصور دا دے
چې که $!A$ د معلومېدو وړ وي، نو په کوم ممکن راتلونکي معرفتي حالت
کښې به~$!A$ معلوم شي۔ څرنګه چې $\lfalse$ معلومولے نۀ شو، نو
په هغه راتلونکي معرفتي حالت کښې به $!A$ معلوم وي خو~$\lfalse$
به معلوم نۀ وي۔

په دې تفسير کښې د منځني حالت د ايستلو اصل نۀ چلېږي۔ ځکه ځينې
$!A$ شته چې لا مو نۀ دي معلوم، خو وروسته ښايي معلوم شي۔ د داسې
!!a{formula}~$!A$ لپاره نۀ $!A$ معلوم دے او نۀ~$\lnot !A$؛
نو $!A \lor \lnot !A$ معلوم نۀ دے۔ خو د بېلګې په توګه موږ
پوهيږو چې $\lnot(!A \land \lnot !A)$ دے۔ ځکه داسې راتلونکے
حالت ممکن نۀ دے چې پکې $!A$ او $\lnot !A$ دواړه معلوم وي؛
دا خبره له دې پرته هم معلومه ده چې~$!A$ يا $\lnot !A$ معلوموو
که نه۔

!!^{relational model}s د شهودي منطق يوازينۍ شونې معنايي پوهنه نۀ ده۔
توپولوژيکي معنايي پوهنه بله لاره ده: دلته قضیې په توپولوژيکي فضا
کښې د پرانيستو سټونو په توګه او نښلوونکي پر دغو سټونو د عملونو په
توګه تفسيرېږي (مثلاً $\land$ له تقاطع سره سمون لري)۔

\end{document}
